\section{Related Work}
\label{sec:related}

\subsection{Harness Engineering and Evaluation for Coding Agents}
\label{ssec:related_harness}
Harness engineering refers to the practice of designing the system surrounding the model, including its tools, interfaces, memory, execution constraints, and feedback loops, which together shape what an agent can do on long-horizon tasks~\citep{rajasekaranHarnessDesignLongrunning2026,lopopoloHarnessEngineeringLeveraging2026,trivedyImprovingDeepAgents2026,anthropicClaudecode2025,steinbergerOpenClawPersonalAI2026,HermesAgentAgent}.
Concretely, the harness mediates how the model perceives and acts on its environment: it exposes the action and observation interfaces over which tool-augmented reasoning unfolds~\citep{anthropicClaudecode2025}, custom agent-computer interfaces for repository navigation, file editing, and command execution~\citep{yangSWEagentAgentComputerInterfaces2024}, as well as sandboxed execution and orchestration support that keep long-horizon runs reproducible~\citep{wangOpenHandsOpenPlatform2025}.

Verifying that such systems actually help has driven the parallel maturation of coding-agent evaluation along two axes: task horizon and environmental realism.
Coverage extends from short-horizon function-level benchmarks focused on contamination and freshness control~\citep{zhuoBigCodeBenchBenchmarkingCode2024,jainLiveCodeBenchHolisticContamination2024}, through repository-scale executable patch resolution~\citep{jimenezSWEbenchCanLanguage2023,yangSWEbenchMultimodalAI2024,dengSWEBenchProCan2025}, to multi-hour, terminal-driven workflows that exercise long-horizon, realistic execution~\citep{miserendinoSWELancerCanFrontier2025,chanMLEbenchEvaluatingMachine2024,merrillTerminalBenchBenchmarkingAgents2026a}.
A parallel infrastructure track packages executable runtimes and verifiers around these benchmarks~\citep{panTrainingSoftwareEngineering2025,jainR2EGymProceduralEnvironment2025,zengSWEHubUnifiedProduction2026}, whose attention to reproducible, traceable, and verifiable execution directly motivates the observation system AHE builds on.


\subsection{Automated Optimization of LLM Agents}
\label{ssec:related_auto}

Approaches to automated agent optimization differ in what evidence the optimizer observes and what it can edit.
Some revise the agent's own outputs through episodic critique and reflection~\citep{madaanSelfRefineIterativeRefinement2023c,shinnReflexionLanguageAgents2023c,guoCritiQMiningData2025}.
Others target prompts and instructions~\citep{khattabDSPyCompilingDeclarative2023}: structured playbooks~\citep{zhangAgenticContextEngineering2025a}, semantic-advantage priors~\citep{caiTrainingFreeGroupRelative2025}, jointly optimized instruction-demonstration pipelines for multi-stage programs~\citep{opsahl-ongOptimizingInstructionsDemonstrations2024a}, and reflective updates driven by Pareto-frontier traces~\citep{agrawalGEPAReflectivePrompt2025a}.
A separate line edits program structure itself, in the form of skill libraries~\citep{wangVoyagerOpenEndedEmbodied2023a}, scored program and agent archives evolved through mutation~\citep{novikovAlphaEvolveCodingAgent2025a,huAutomatedDesignAgentic2024}, and graph-structured workflows searched or learned from rollouts~\citep{zhangAFlowAutomatingAgentic2024,zhouSymbolicLearningEnables2024a}.

AHE tunes the full harness as a combinatorial whole rather than a single editable surface, so cross-component trade-offs become legible to the optimizer.
It also keeps the human prior minimal, leaving methodology for the optimizer to discover from rollouts rather than fixing it by hand.
% The optimization target is long-horizon software-engineering work.
We describe the substrate, trajectory analysis, and iteration that realize these choices in Section~\ref{sec:method}.
